\subsection{From categorical modes to complete replayed responses}
\label{app:theory-exemplar-bridge}

The categorical analysis protects an execution key, whereas the learner
scores a particular response string. The bridge is event containment: one
complete verified response is one way to produce its key. A bound on its
sequence surprisal therefore protects the key even when many other strings
realize the same behavior. Length normalization changes the numerical loss
scale, so its factor must remain explicit in that bound.

\begin{lemma}[Retention from a joint exemplar potential]
\label{lem:shared-exemplar-retention}
Fix finitely many verified complete-response exemplars $(x_j,e_j)$, each
realizing key $b_j$, with lengths $L_j>0$ and fixed weights $\alpha_j>0$.
Let
\[
 F(\theta)=\sum_j\alpha_j
 \frac{-\log\pi_\theta(e_j\mid x_j)}{L_j}-J(\theta),
 \qquad J(\theta)\le J_{\max}<\infty.
\]
Suppose $F(\theta(T))$ is finite and
$F(\theta(t))\le F(\theta(T))$ for every $t\ge T$. With
$D=F(\theta(T))+J_{\max}$, every protected exemplar satisfies
\[
 p_\theta(b_j\mid x_j)\ge\pi_\theta(e_j\mid x_j)
 \ge\exp\!\left(-\frac{L_jD}{\alpha_j}\right)>0
 \qquad(t\ge T).
\]
\end{lemma}
\begin{proof}
The weighted exemplar loss equals $F(\theta(t))+J(\theta(t))\le D$.
Each summand is nonnegative, so
$-\log\pi_\theta(e_j\mid x_j)\le L_jD/\alpha_j$.
The complete-response event is a subset of the event producing key $b_j$,
which proves the stated lower bound on every protected key's probability.
\end{proof}

\paragraph{Example: a token average is not a response probability.}
Suppose a complete response has four scored tokens, including termination,
and each token has conditional probability $1/2$ along that response. Its
mean log likelihood is $-\log2$, but its complete-response probability is
$(1/2)^4=1/16$; exponentiating the mean gives $1/2$, the geometric mean of
token probabilities. For a single protected exemplar with $L=4$, $\alpha=1$,
and joint energy ceiling $D=\log2$, the lemma gives exactly the $1/16$ floor.
If a different complete response with the same key has probability $1/32$,
the two disjoint response events together give key probability at least
$1/16+1/32=3/32$. The exemplar bound can thus be conservative even when its
sequence probability is known exactly.

Exact gradient flow of a smooth joint potential is one sufficient condition
for this monotonicity; the lemma requires only the stated energy inequality.
Shared parameters and multiple prompts therefore need not invalidate
retention. For example, the exact categorical mean model admits
$J(\theta)=\sum_x\nu_x\Psi_x(P_x(\theta))$ with fixed nonnegative prompt
weights and $J_{\max}=\sum_x\nu_x\Psi_x(1)$. The bound does not require
uniform key weights or imply convergence to a uniform policy. It does not
establish joint energy descent for the implemented stochastic, clipped,
alternating optimizer. The likelihood must describe the complete response
under the sampling law being bounded, including its termination event or a
fixed deterministic horizon. A prefix likelihood alone, or a score under a
different temperature or token-truncation rule, is insufficient.

\paragraph{Retention and finite evaluation visibility.}
At a fixed checkpoint, suppose $k$ distinct banked keys for one prompt have
probability floors $p_b\ge\delta_b>0$. For $K$ independent draws from that
same policy, the banked distinct count $D_{\mathcal B,K}$ satisfies
\[
 \begin{aligned}
 \mathbb E[D_{\mathcal B,K}]
 &\ge\sum_{b\in\mathcal B}\left[1-(1-\delta_b)^K\right],\\
 \Pr(\text{some banked key is unseen})
 &\le\sum_{b\in\mathcal B}(1-\delta_b)^K
 \le k e^{-K\delta_{\min}},
 \end{aligned}
\]
where $\delta_{\min}=\min_b\delta_b$. The first statement follows from
indicator expectations and the second from the union bound. Thus
$K\ge\log(k/\varepsilon)/\delta_{\min}$ suffices to observe all banked keys
with probability at least $1-\varepsilon$, for $0<\varepsilon<1$.
The energy floors can make this budget impractically large; asymptotic
retention does not imply visibility in an eight-draw evaluation.

\paragraph{Example: protection can still be hard to observe.}
A key with probability exactly $1/16$ appears in eight independent draws
with probability $1-(15/16)^8\approx0.403$. If four keys each have a $1/16$
floor, the sufficient budget above is
$\lceil16\log(4/0.05)\rceil=71$ draws for at least $95\%$ probability of
seeing all four. These hypothetical probabilities distinguish retaining a
mode from reliably observing it in a short evaluation.

